\documentclass[10pt,twocolumn]{article}
\usepackage[margin=0.48in]{geometry}
\usepackage{amsmath,amssymb,mathtools,bm}
\usepackage{graphicx}
\usepackage[most]{tcolorbox}
\usepackage{booktabs}
\usepackage{enumitem}
\usepackage{xcolor}
\usepackage{hyperref}
\usepackage{caption}
\usepackage{microtype}
\hypersetup{colorlinks=true,linkcolor=blue!45!black,urlcolor=blue!45!black,citecolor=blue!45!black}
\setlength{\columnsep}{0.18in}
\setlength{\parindent}{0pt}
\setlength{\parskip}{2.2pt}
\setlist[itemize]{leftmargin=1.1em,itemsep=1pt,topsep=2pt}
\setlist[enumerate]{leftmargin=1.3em,itemsep=1pt,topsep=2pt}
\captionsetup{font=footnotesize,skip=3pt}
\allowdisplaybreaks
\newcommand{\R}{\mathbb R}
\newcommand{\range}{\operatorname{range}}
\newcommand{\diag}{\operatorname{diag}}
\newcommand{\spanop}{\operatorname{span}}
\newcommand{\nov}{\mathcal N}
\newcommand{\eps}{\varepsilon}
\newcommand{\ket}[1]{\lvert #1\rangle}
\newcommand{\bra}[1]{\langle #1\rvert}
\newcommand{\braket}[2]{\langle #1\mid #2\rangle}
\definecolor{deepblue}{HTML}{1E5575}
\definecolor{pale}{HTML}{EFF7FA}
\definecolor{warm}{HTML}{FFF4E8}
\definecolor{greenbox}{HTML}{EEF8EF}
\definecolor{redbox}{HTML}{FFF0F0}
\tcbset{boxsep=2pt,left=4pt,right=4pt,top=3pt,bottom=3pt,arc=2pt,boxrule=0.55pt}
\newtcolorbox{intuitionbox}{colback=pale,colframe=deepblue,title=Intuition,fonttitle=\bfseries\small}
\newtcolorbox{jargonbox}{colback=warm,colframe=orange!60!black,title=Jargon-heavy statement,fonttitle=\bfseries\small}
\newtcolorbox{cautionbox}{colback=redbox,colframe=red!55!black,title=Caution,fonttitle=\bfseries\small}

\title{Paper-II McLachlan Projection Geometry Primer\\[-1mm]
\large Same Gram/Schur machinery, written only in time-dynamics notation}
\author{Local pedagogical note}
\date{May 17, 2026}
\begin{document}
\maketitle

\begin{abstract}
This note rewrites the nonorthogonal projection primer entirely in Paper-II McLachlan time-dynamics notation. The recurring object is a horizontal tangent block \(\bar T_k\) and a horizontal Schr\"odinger drift \(\bar b_k\). McLachlan dynamics is the least-squares projection of \(\bar b_k\) onto \(\range(\bar T_k)\):
\[
K_k\dot\theta_k=f_k,
\qquad
\epsilon_{{\rm expr},k}^2=\|\bar b_k\|^2-f_k^\top K_k^+f_k.
\]
Append-block Schur gain is the same residual-projection calculation after splitting old tangent coordinates from candidate append coordinates.
\end{abstract}

\section{The McLachlan projection question}

\begin{intuitionbox}
At time step \(k\), the exact Schr\"odinger drift points in some horizontal Hilbert-space direction. The ansatz can move only through its horizontal tangent columns. McLachlan chooses velocity coefficients \(\dot\theta_k\) so the synthesized tangent velocity \(\bar T_k\dot\theta_k\) best matches the drift \(\bar b_k\).
\end{intuitionbox}

\begin{jargonbox}
Paper II is least-squares vector-field projection on a time-dependent variational tangent bundle. The horizontal tangent block \(\bar T_k\) is the current nonorthogonal dictionary, the Schr\"odinger drift \(\bar b_k\) is the target vector field, and the McLachlan normal equations convert raw force overlaps into velocity coordinates:
\[
K_k\dot\theta_k=f_k.
\]
The residual scalar \(\|\bar b_k\|^2-f_k^\top K_k^+f_k\) is the drift norm not represented by the current ansatz tangent space.
\end{jargonbox}

Let
\begin{equation}
\bar T_k=[\bar\tau_{1,k}\ \cdots\ \bar\tau_{N_k,k}],
\qquad
\bar b_k=-i(H_k-\langle H_k\rangle_{\psi_k})\ket{\psi_k}.
\label{eq:p2_Tb}
\end{equation}
The variational velocity is synthesized as
\begin{equation}
\bar T_k\dot\theta_k=\sum_i\dot\theta_{k,i}\bar\tau_{i,k}.
\label{eq:p2_velocity_synthesis}
\end{equation}
Choose \(\dot\theta_k\) by least squares:
\begin{equation}
\dot\theta_k^\star\in\arg\min_v\|\bar b_k-\bar T_kv\|^2.
\label{eq:p2_ls}
\end{equation}
Expanding gives
\begin{align}
\|\bar b_k-\bar T_kv\|^2
&=\langle \bar b_k-\bar T_kv,\bar b_k-\bar T_kv\rangle
\notag\\
&=\|\bar b_k\|^2
-2v^\top\Re(\bar T_k^\dagger\bar b_k)
+v^\top\Re(\bar T_k^\dagger\bar T_k)v
\notag\\
&=\|\bar b_k\|^2-2v^\top f_k+v^\top G_kv,
\label{eq:p2_expand}\\
G_k&=\Re(\bar T_k^\dagger\bar T_k),
\qquad
f_k=\Re(\bar T_k^\dagger\bar b_k).
\label{eq:p2_Gf}
\end{align}
With regularization \(\Lambda_k\), define
\begin{equation}
K_k=G_k+\Lambda_k.
\label{eq:p2_K}
\end{equation}
The normal equations are
\begin{equation}
\boxed{K_k\dot\theta_k=f_k.}
\label{eq:p2_normal}
\end{equation}
Here \(f_k\) is the force-overlap vector: it records how strongly the target drift shadows each available tangent direction. The solve through \(K_k\) corrects those raw shadows for nonorthogonality and regularization before returning velocity coordinates.
The projected drift is \(\bar T_k\dot\theta_k\), and the expressivity miss is
\begin{equation}
\boxed{
\epsilon_{{\rm expr},k}^2
=\|\bar b_k\|^2-f_k^\top K_k^+f_k.
}
\label{eq:p2_expr_miss}
\end{equation}

\section{Analysis, synthesis, and the residual scalar}

The McLachlan analysis map is
\begin{equation}
\bar b_k\longmapsto f_k=\Re(\bar T_k^\dagger\bar b_k),
\label{eq:p2_analysis}
\end{equation}
and the synthesis map is
\begin{equation}
v\longmapsto\bar T_kv.
\label{eq:p2_synthesis_map}
\end{equation}
The normal solve produces the best velocity coordinates:
\begin{equation}
\begin{aligned}
K_kv^\star=f_k
&\Longrightarrow v^\star=K_k^+f_k\\
&\Longrightarrow \bar T_kv^\star
\text{ is the fitted variational drift}\\
&\Longrightarrow
\|\bar b_k-\bar T_kv^\star\|^2
=\|\bar b_k\|^2-f_k^\top K_k^+f_k.
\end{aligned}
\label{eq:p2_chain}
\end{equation}
Thus the miss \(\epsilon_{{\rm expr},k}^2\) is the dynamical analogue of Paper-I novelty: it is the part of the target vector field that survives projection against the current nonorthogonal tangent span.

\section{Supported tangent geometry}

Diagonalize the McLachlan normal metric:
\begin{equation}
K_k=V_k\diag(s_{k,1},\ldots,s_{k,N_k})V_k^\top.
\label{eq:p2_eigs}
\end{equation}
The columns \(v_{k,a}\) are obtained from
\begin{equation}
K_kv_{k,a}=s_{k,a}v_{k,a},
\qquad \|v_{k,a}\|=1.
\label{eq:p2_eigenproblem}
\end{equation}
Each \(v_{k,a}\) is a velocity-coordinate recipe:
\begin{equation}
\bar T_kv_{k,a}=\sum_i(v_{k,a})_i\bar\tau_{i,k}.
\label{eq:p2_mode_synthesis}
\end{equation}
Its realized tangent-velocity length is
\begin{equation}
\|\bar T_kv_{k,a}\|^2
=v_{k,a}^\top G_kv_{k,a},
\label{eq:p2_mode_length_G}
\end{equation}
and, without regularization, this is the corresponding eigenvalue of \(G_k\). Small supported modes are velocity mixtures that barely move the state but can receive large fitted coefficients.

A supported inverse uses the resolved metric modes:
\begin{equation}
K_{k,\tau}^+
=V_k\diag\left(\frac{\mathbf 1[s_{k,a}\ge\tau s_{k,\max}]}{s_{k,a}}\right)V_k^\top.
\label{eq:p2_supported_inverse}
\end{equation}

\section{Append-block Schur gain}

Split the tangent columns into current columns \(\bar T_k\) and candidate append columns \(\bar U_{r,k}\). Define
\begin{align}
K_k&=\Re(\bar T_k^\dagger\bar T_k)+\Lambda_k,
&
B_{r,k}&=\Re(\bar T_k^\dagger\bar U_{r,k}),
\notag\\
C_{r,k}&=\Re(\bar U_{r,k}^\dagger\bar U_{r,k})+\Lambda_{r,k},
&
q_{r,k}&=\Re(\bar U_{r,k}^\dagger\bar b_k),
\notag\\
f_k&=\Re(\bar T_k^\dagger\bar b_k).
\label{eq:p2_append_blocks}
\end{align}
The augmented normal matrix is
\begin{equation}
\begin{pmatrix}
K_k&B_{r,k}\\
B_{r,k}^\top&C_{r,k}
\end{pmatrix}.
\label{eq:p2_augmented_normal}
\end{equation}
The augmented force vector is
\begin{equation}
\begin{pmatrix}
f_k\\ q_{r,k}
\end{pmatrix}.
\label{eq:p2_augmented_force}
\end{equation}
Eliminating old velocity coordinates gives the candidate residual geometry
\begin{equation}
S_{r,k}=C_{r,k}-B_{r,k}^\top K_k^+B_{r,k},
\label{eq:p2_append_schur}
\end{equation}
and the residual force seen by the candidate block
\begin{equation}
w_{r,k}=q_{r,k}-B_{r,k}^\top K_k^+f_k.
\label{eq:p2_residual_force}
\end{equation}
The append gain is
\begin{equation}
\boxed{
\Delta_{r,k}^{\rm Schur}=w_{r,k}^\top S_{r,k}^+w_{r,k}.
}
\label{eq:p2_append_gain}
\end{equation}
This is the amount of remaining drift the candidate block can explain after the current tangent coordinates have already had their best chance.

\begin{jargonbox}
The Schur block \(S_{r,k}\) is the candidate append metric after quotienting out old tangent velocities. The vector \(w_{r,k}\) is the residual force seen by the append block after the old tangent coordinates have already explained their share of \(\bar b_k\).
\end{jargonbox}

\section{Stay, append, prune as projection decisions}

The same scalar geometry supports three actions:
\begin{equation}
\begin{array}{ccl}
\text{stay}&:&\epsilon_{{\rm expr},k}^2\text{ is already small},\\[1mm]
\text{append}&:&\Delta_{r,k}^{\rm Schur}\text{ is large for some }r,\\[1mm]
\text{prune}&:&\text{a retained tangent mode has negligible supported contribution}.
\end{array}
\label{eq:p2_actions}
\end{equation}
All three decisions are metric-projection decisions inside the same nonorthogonal tangent bundle.

\appendix
\section{Tiny McLachlan example}

Let
\begin{equation}
\bar T=\begin{pmatrix}1&1\\0&1\\0&0\end{pmatrix},
\qquad
\bar b=\begin{pmatrix}0\\1\\1\end{pmatrix}.
\end{equation}
Then
\begin{equation}
G=\bar T^\top\bar T=\begin{pmatrix}1&1\\1&2\end{pmatrix},
\qquad
f=\bar T^\top\bar b=\begin{pmatrix}0\\1\end{pmatrix}.
\end{equation}
Solving \(G\dot\theta=f\) gives \(\dot\theta=(-1,1)^\top\), so
\begin{equation}
\bar T\dot\theta=\begin{pmatrix}0\\1\\0\end{pmatrix}.
\end{equation}
The third coordinate of \(\bar b\) remains outside the current tangent span.

\section{Example: force overlaps are not velocity coefficients}

Let
\begin{equation}
\bar T=
\begin{pmatrix}
1&1\\
0&1
\end{pmatrix},
\qquad
\bar b=
\begin{pmatrix}
1\\0
\end{pmatrix}.
\end{equation}
Then
\begin{equation}
G=\bar T^\top\bar T=
\begin{pmatrix}1&1\\1&2\end{pmatrix},
\qquad
f=\bar T^\top\bar b=
\begin{pmatrix}1\\1\end{pmatrix}.
\end{equation}
Solving \(Gv=f\) gives
\begin{equation}
v^\star=\begin{pmatrix}1\\0\end{pmatrix}.
\end{equation}
The force overlaps \(f=(1,1)^\top\) are both large because both tangent columns shadow \(\bar b\). The velocity coefficients are \((1,0)^\top\) after the Gram solve removes the double-counting.

\section{Example: append-block Schur gain}

Let
\begin{equation}
K=(2),\qquad
B=(1),\qquad
C=(3),\qquad
f=(4),\qquad
q=(5).
\end{equation}
Then
\begin{equation}
S=C-B^\top K^{-1}B=3-\frac12=\frac52,
\qquad
w=q-B^\top K^{-1}f=5-2=3.
\end{equation}
The append gain is
\begin{equation}
\Delta^{\rm Schur}=w^\top S^{-1}w
=3\left(\frac{2}{5}\right)3
=\frac{18}{5}.
\end{equation}
This number is the residual drift explained by the candidate after the old tangent coordinate has been eliminated.

\end{document}
